% -------------------- Como usar este livro --------------------------
\chapter*{Como Usar Este Livro}
\addcontentsline{toc}{chapter}{Como Usar Este Livro}
\markboth{Como Usar Este Livro}{Como Usar Este Livro}

Este livro foi projetado para ser lido \emph{com o computador ligado}. Cada
trecho de código é curto o bastante para ser digitado à mão --- e recomendamos
que você o digite, pois o aprendizado se consolida pelos dedos, não apenas pelos
olhos. Ao longo do texto você encontrará elementos visuais recorrentes:

\begin{caixanota}
Aprofunda um conceito, esclarece uma sutileza ou conecta o assunto a
capítulos futuros. Ler é recomendável, mas o fluxo principal não depende disso.
\end{caixanota}

\begin{caixadica}
Truques práticos, atalhos e boas práticas que economizam tempo e evitam
frustrações comuns.
\end{caixadica}

\begin{caixaatencao}
Erros frequentes e armadilhas. Preste atenção especial: quase todo iniciante
tropeça aqui pelo menos uma vez.
\end{caixaatencao}

\begin{caixamaos}
Exercícios guiados de digitar-e-executar. É onde a teoria vira prática.
Não pule estas seções.
\end{caixamaos}

\begin{caixaancora}
Sinaliza um avanço no \textbf{Projeto Âncora}, o fio condutor que atravessa todo
o livro. Cada etapa acrescenta uma peça ao seu \emph{pipeline} de simulação de
\qubit{}s.
\end{caixaancora}

\noindent\textbf{Convenções tipográficas.} Nomes de comandos, funções e trechos
curtos de código aparecem em \code{fonte monoespaçada}. Blocos de código \python{}
são numerados à esquerda e destacados em fundo cinza. Saídas de terminal
aparecem em fundo escuro.

\begin{caixadica}
Todo o código deste livro estará disponível em um repositório on-line, em forma
de \emph{notebooks} Jupyter prontos para executar. Ainda assim, digite os
exemplos ao menos uma vez: a fluência não se baixa, se constrói.
\end{caixadica}
